\subsection{RQ3: Admission Decision Quality}
\label{sec:rq3}

RQ3 evaluates admission decisions rather than deployed stage execution.
We audited 3{,}746 decisions on the S26 Ultra under an always-admit configuration.
Optional stages always executed with the watchdog suppressed; admission decisions were recorded but not enforced.
Each decision is thus scored against its realized outcome: feasible if its remaining sequence completed within the remaining timeout budget at that decision, and unsafe otherwise.
Because no decision was enforced, every approach is scored on the same decisions and outcomes.

\smallskip\noindent\textbf{Comparison with naive approaches.}
Table~\ref{tab:rq3_naive} shows that our approach admits far fewer unsafe decisions than naive approaches.
These approaches replace \(U\) with the mean or the maximum of each remaining stage's last three observed durations, or with the point estimate \(\hat P\) alone.
Our approach admitted only 5 of the 178 unsafe decisions, whereas the mean, the maximum, and \(\hat P\) admitted 42, 14, and 47, respectively.
In exchange, it skipped 2.7\% of feasible decisions, against 0.1--1.8\%.
In particular, \(\hat P\) shares our stage-duration model, so residual calibration alone rejected 42 of the 47 unsafe decisions that \(\hat P\) admitted.
Nor can the safeguards absorb this gap: the watchdog bounds only the admitted stage, leaving later overruns to later skips made with the same naive estimate, and its budget in Equation~\ref{eq:watchdog} itself reserves the calibrated \(U(\mathcal E_i)\).
Estimating stage durations is therefore not enough; admission must also account for how far those estimates have been exceeded.

\smallskip\noindent\textbf{Decision accuracy.}
Admission was accurate on both sides of the budget boundary: it recommended execution for 97.3\% of feasible decisions and skipping for 97.2\% of unsafe ones.
Table~\ref{tab:rq3_admission_audit} shows that this accuracy held in every condition--stage group, at 94.7--99.4\% and 94.3--100\%, respectively.
Its conservative skips, 2.7\% of feasible decisions, were marginal cases whose median remaining-sequence margins were only 0.7--2.6\% of the timeout budget across the condition--stage groups.
Its correct skips caught remaining sequences that exceeded the budget by a median of only 3.0--3.5\%.

\smallskip\noindent\textbf{Unsafe admits and safeguards.}
Five unsafe decisions (2.8\%) received admit recommendations.
Figure~\ref{fig:rq3_unsafe_spike_anatomy} compares the remaining-sequence latency and CPU time of these five decisions with those of the preceding capture.
M1--M3 and S1--S2 label the multi-frame and single-frame decisions, respectively.
For M1 and M2, latency rose to 2.32\(\times\) and 1.95\(\times\), whereas CPU time rose only to 1.41\(\times\) and 1.15\(\times\).
The busy-core proxy, defined as CPU time per unit of wall time, consequently fell to 0.60\(\times\) and 0.59\(\times\).
These misses coincided with latency increases that were not captured by the decision-time estimates.

In trace replay, the layered safeguards covered all five admission misses.
The replay re-applied the safeguard logic to the recorded stage durations, budgets, upper estimates, and admission decisions.
For M1 and M2, the admitted stage exceeded its watchdog budget, 1{,}460~ms against 1{,}202~ms and 1{,}648~ms against 889~ms, respectively, so fallback would trigger before the stage returned.
For M3, a later decision would skip \(S\), reducing the remaining processing time.
For S1 and S2, an earlier decision would skip \(M\), leaving more budget when execution reached \(S\).
With these skips applied to the recorded durations, the remaining sequence fit within the adjusted budget in all three cases.

\begin{table}[t]
  \centering
  \setlength{\abovecaptionskip}{2pt}%
  \setlength{\belowcaptionskip}{3pt}%
  \caption{RQ3: Comparison with naive approaches.}
  \label{tab:rq3_naive}

  {\scriptsize
  \renewcommand{\arraystretch}{1.12}%
  \newlength{\rqnfawidth}\settowidth{\rqnfawidth}{99.9\% (3{,}565)}%
  \newlength{\rqnfswidth}\settowidth{\rqnfswidth}{2.7\% (98)}%
  \newlength{\rqnuawidth}\settowidth{\rqnuawidth}{26.4\% (47)}%
  \newlength{\rqnuswidth}\settowidth{\rqnuswidth}{97.2\% (173)}%
  \newcommand{\rqnfa}[1]{\makebox[\rqnfawidth][r]{#1}}%
  \newcommand{\rqnfs}[1]{\makebox[\rqnfswidth][r]{#1}}%
  \newcommand{\rqnua}[1]{\makebox[\rqnuawidth][r]{#1}}%
  \newcommand{\rqnus}[1]{\makebox[\rqnuswidth][r]{#1}}%
  \newsavebox{\rqnbox}%

  \newcommand{\rqntabular}{%
  \begin{tabular}{@{}%
      >{\raggedright\arraybackslash}p{78pt}|%
      cc|cc%
    @{\hspace{\tabcolsep}}}
    \toprule
    \multirow[c]{3}{=}[-1.6pt]{\centering\textbf{Approach}}
      & \multicolumn{4}{c}{\textbf{Admission audit, \% (\(n\))}} \\
    \cmidrule(l){2-5}
      & \multicolumn{2}{c|}{\textbf{Feasible decisions}}
      & \multicolumn{2}{c}{\textbf{Unsafe decisions}} \\
    \cmidrule(lr){2-3}\cmidrule(l){4-5}
      & \multicolumn{1}{c}{\textbf{Admitted}}
      & \multicolumn{1}{c|}{\textbf{Skipped}}
      & \multicolumn{1}{c}{\textbf{Admitted}}
      & \multicolumn{1}{c}{\textbf{Skipped}} \\
    \midrule

    \multicolumn{5}{@{}l}{\textbf{Naive approaches}} \\
    \addlinespace[1pt]
    Mean of last 3 stage times & \rqnfa{99.6\% (3{,}554)} & \rqnfs{0.4\% (14)} & \rqnua{23.6\% (42)} & \rqnus{76.4\% (136)} \\
    Max of last 3 stage times  & \rqnfa{98.2\% (3{,}505)} & \rqnfs{1.8\% (63)} & \rqnua{7.9\% (14)} & \rqnus{92.1\% (164)} \\
    Point estimate \(\hat P\)  & \rqnfa{99.9\% (3{,}565)} & \rqnfs{0.1\% (3)} & \rqnua{26.4\% (47)} & \rqnus{73.6\% (131)} \\
    \midrule

    \multicolumn{5}{@{}l}{\textbf{Our approach}} \\
    \addlinespace[1pt]
    Residual calibration \(U\) & \rqnfa{97.3\% (3{,}470)} & \rqnfs{2.7\% (98)} & \rqnua{2.8\% (5)} & \rqnus{97.2\% (173)} \\
    \bottomrule
  \end{tabular}}%
  \fittabcolsep{\rqnbox}{\rqntabular}{\columnwidth}{10}%
  \rqntabular
  }
\end{table}

% Notes: docs/exhibits.md#tab_rq3_admission_quality
\begin{table}[t]
  \centering
  \setlength{\abovecaptionskip}{2pt}%
  \setlength{\belowcaptionskip}{3pt}%
  \caption{RQ3: Admission decision quality.}
  \label{tab:rq3_admission_audit}

  {\scriptsize
  \renewcommand{\arraystretch}{1.12}%
  \newlength{\rqauditwidth}%
  \settowidth{\rqauditwidth}{97.3\% (3{,}470)}%
  \newcommand{\rqauditcell}[1]{\makebox[\rqauditwidth][r]{#1}}%
  \newsavebox{\rqtwobox}%

  \newcommand{\rqtwotabular}{%
  \begin{tabular}{@{}%
      >{\raggedright\arraybackslash}p{50pt}|%
      cc|cc%
    @{\hspace{\tabcolsep}}}
    \toprule
    \multirow[c]{3}{=}[-1.6pt]{\centering\textbf{Group}}
      & \multicolumn{4}{c}{\textbf{Admission audit, \% (\(n\))}} \\
    \cmidrule(l){2-5}
      & \multicolumn{2}{c|}{\textbf{Feasible decisions}}
      & \multicolumn{2}{c}{\textbf{Unsafe decisions}} \\
    \cmidrule(lr){2-3}\cmidrule(l){4-5}
      & \multicolumn{1}{c}{\textbf{Admitted}}
      & \multicolumn{1}{c|}{\textbf{Skipped}}
      & \multicolumn{1}{c}{\textbf{Admitted}}
      & \multicolumn{1}{c}{\textbf{Skipped}} \\
    \midrule

    \multicolumn{5}{@{}l}{\textbf{12MP, normal}} \\
    \addlinespace[1pt]
    Multi-frame  & \rqauditcell{98.1\% (831)} & \rqauditcell{1.9\% (16)} & \rqauditcell{\textbf{5.7\% (2)}} & \rqauditcell{94.3\% (33)} \\
    Single-frame & \rqauditcell{99.4\% (842)} & \rqauditcell{0.6\% (5)} & \rqauditcell{\textbf{0.0\% (0)}} & \rqauditcell{100.0\% (35)} \\
    \midrule

    \multicolumn{5}{@{}l}{\textbf{24MP, mem.\ pressure}} \\
    \addlinespace[1pt]
    Multi-frame  & \rqauditcell{94.7\% (892)} & \rqauditcell{5.3\% (50)} & \rqauditcell{\textbf{1.9\% (1)}} & \rqauditcell{98.1\% (53)} \\
    Single-frame & \rqauditcell{97.1\% (905)} & \rqauditcell{2.9\% (27)} & \rqauditcell{\textbf{3.7\% (2)}} & \rqauditcell{96.3\% (52)} \\
    \midrule
    \textbf{Overall} & \rqauditcell{97.3\% (3{,}470)} & \rqauditcell{2.7\% (98)} & \rqauditcell{\textbf{2.8\% (5)}} & \rqauditcell{97.2\% (173)} \\
    \bottomrule
  \end{tabular}}%
  \fittabcolsep{\rqtwobox}{\rqtwotabular}{\columnwidth}{10}%
  \rqtwotabular
  }
\end{table}

% Notes: docs/exhibits.md#fig_rq3_unsafe_spike_anatomy
\begin{figure}[t]
  \centering
  \setlength{\abovecaptionskip}{3pt}%
  \setlength{\belowcaptionskip}{2pt}%
  \newcommand{\rqtwosg}[2]{%
    \node[font=\tiny,inner sep=0pt,anchor=west] at (axis cs:1.66,#1)
      {\ding{51}\,\textsf{#2}};}%
  \newcommand{\rqtwosgtwo}[3]{%
    \node[font=\tiny,inner sep=0pt,anchor=west] at (axis cs:1.66,#1)
      {\begin{tabular}[c]{@{}c@{\hspace{2pt}}l@{}}
        \multirow{2}{*}{\ding{51}} & \textsf{#2} \\[-1pt]
        & \textsf{#3}
      \end{tabular}};}%
  \begin{tikzpicture}
    \newcommand{\rqtwovalr}[4]{%
      \node[font=\tiny,inner sep=0pt,anchor=west,yshift=#1,xshift=2.5pt]
        at (axis cs:#2,#3) {$\times$#4};}
    \newcommand{\rqtwovall}[4]{%
      \node[font=\tiny,inner sep=0pt,anchor=east,yshift=#1,xshift=-2.5pt]
        at (axis cs:#2,#3) {$\times$#4};}
    \begin{axis}[
      xbar,
      bar width=6.5pt,
      y=25pt,
      scale only axis=true,
      width=0.720\columnwidth,
      ymin=0.4, ymax=5.6,
      ytick={1,2,3,4,5},
      yticklabels={S2,S1,M3,M2,M1},
      xmin=-0.65, xmax=1.60,
      xtick={-0.4,-0.2,0,0.2,0.4,0.6,0.8,1.0,1.2},
      xticklabels={0.6,0.8,1.0,1.2,1.4,1.6,1.8,2.0,2.2},
      xlabel={Ratio to the previous capture ($\times$)},
      axis x line*=bottom,
      axis y line*=left,
      tick align=outside,
      xmajorgrids,
      grid style={black!14,line width=0.3pt},
      axis line style={line width=0.35pt},
      tick style={line width=0.35pt},
      tick label style={font=\tiny},
      label style={font=\tiny},
      xlabel style={yshift=2pt},
      legend style={
        font=\tiny,
        draw=black!25,
        fill=white,
        line width=0.2pt,
        inner sep=2pt,
        row sep=-1pt,
        at={(axis cs:1.56,1.5)},
        anchor=east,
        legend columns=1,
        /tikz/every even column/.append style={column sep=3pt},
      },
      legend cell align=left,
      legend image code/.code={%
        \draw[#1] (0cm,-0.075cm) rectangle (0.24cm,0.075cm);},
      clip=false,
    ]

    \draw[black!60,densely dashed,line width=0.5pt]
      (axis cs:0,0.4) -- (axis cs:0,5.6);
    \draw[black!70,line width=0.6pt]
      (axis cs:-0.65,2.5) -- (axis cs:1.60,2.5);

    \addplot[draw=black,fill=black,line width=0.3pt,bar shift=6.5pt]
      coordinates {(1.32,5) (0.95,4) (0.14,3)
                   (0.60,2) (0.42,1)};
    \addlegendentry{Latency}

    \addplot[draw=black,fill=black!35,line width=0.3pt,bar shift=0pt]
      coordinates {(0.41,5) (0.15,4) (0.23,3)
                   (0.24,2) (0.37,1)};
    \addlegendentry{CPU time}

    \addplot[draw=black,fill=white,line width=0.35pt,bar shift=-6.5pt]
      coordinates {(-0.40,5) (-0.41,4) (0.08,3)
                   (-0.24,2) (-0.04,1)};
    \addlegendentry{Busy cores}

    \rqtwovalr{6.5pt}{1.32}{5}{2.32}
    \rqtwovalr{6.5pt}{0.95}{4}{1.95}
    \rqtwovalr{6.5pt}{0.14}{3}{1.14}
    \rqtwovalr{6.5pt}{0.60}{2}{1.60}
    \rqtwovalr{6.5pt}{0.42}{1}{1.42}

    \rqtwovalr{0pt}{0.41}{5}{1.41}
    \rqtwovalr{0pt}{0.15}{4}{1.15}
    \rqtwovalr{0pt}{0.23}{3}{1.23}
    \rqtwovalr{0pt}{0.24}{2}{1.24}
    \rqtwovalr{0pt}{0.37}{1}{1.37}

    \rqtwovall{-6.5pt}{-0.40}{5}{0.60}
    \rqtwovall{-6.5pt}{-0.41}{4}{0.59}
    \rqtwovalr{-6.5pt}{0.08}{3}{1.08}
    \rqtwovall{-6.5pt}{-0.24}{2}{0.76}
    \rqtwovall{-6.5pt}{-0.04}{1}{0.96}

    \node[font=\tiny,inner sep=0pt,anchor=south west]
      at (axis cs:1.66,5.62) {\textsf{prevented by}};
    \rqtwosg{5}{watchdog}
    \rqtwosg{4}{watchdog}
    \rqtwosgtwo{3}{later single-frame}{skip}
    \rqtwosgtwo{2}{earlier multi-frame}{skip}
    \rqtwosgtwo{1}{earlier multi-frame}{skip}

    \end{axis}
  \end{tikzpicture}
   \caption{Reason about each unsafe admit of Table~\ref{tab:rq3_admission_audit}}
  \label{fig:rq3_unsafe_spike_anatomy}
\end{figure}

